% Procedencia HMT: output/EXTREMA_MEDIA_RAZON_NARRACION_Y_REVISION_20260916/
% ARTICULO/ES/source/libro/01_acoplamiento.tex, lib01:cadena-registro y lib01:K.
% Dirección: ARTICULO/ES/source/sections/k_direccion_dimensional.tex.
% Estatuto de las redes, cartas y flujos de este apéndice: FORMALIZACION_NUEVA.
% La extracción de K y el proyector P_3 son RESULTADO_RECUPERADO.
\section{Mesure planaire de frontière, coordonnées de Plücker et mutations positives}
\label{app:network-cluster}

La représentation dodécaphasique de l'état enrichi admet une réalisation positive dont le réseau, les coordonnées et les changements de coordinatisation peuvent être écrits intégralement. La composition de départ est
\[
 \mathrm{APP}\longrightarrow\mathrm{TRIT}\longrightarrow\mathrm{TPK}
 \longrightarrow x_{\rm term}\longrightarrow\mathcal R_{12}
 \longrightarrow(b^{90},b^{120},Q_{\rm TPK})
 \longrightarrow U_{\rm sgn}\longrightarrow K.
\]
Cette lecture agit sur l’état enrichi qui conserve les deux feuillets,
l’orientation, la retenue, la mémoire et la frontière. Le registre déjà engendré
est
\[
 K=(234,543,140,729,659,824,621,58,914,794,146,601),
 \qquad Q=\sum_{r=1}^{12}K_r=6263.
\]
Son extraction se récupère exactement à partir des deux canaux et de la charge totale :
\[
 b^{90}=(-495,-116,-684,108,601,-90,-173,-88,313,560,-397,461),
\]
\[
 b^{120}=(-425,-281,-481,671,-255,30,475,-543,680,251,6,-128).
\]
Sur les orbites de positions \((r,r+3,r+6,r+9)\), pour \(r=1,2,3\),
soient \(B_r=\sum_{j=0}^3b^{120}_{r+3j}\) et
\(d_{r,j}=b^{90}_{r+3j}\). On obtient
\[
 s_1=(Q+2B_1+B_2)/3,\quad s_2=s_1-B_1,\quad s_3=s_2-B_2,
\]
\[
 U^{(r)}=(s_r,d_{r,1}-d_{r,3},d_{r,0}+d_{r,2},d_{r,0}-d_{r,2}),
 \qquad K|_r=\tfrac14H_4U^{(r)},
\]
\[
 H_4=\begin{pmatrix}1&1&1&1\\1&1&-1&-1\\1&-1&1&-1\\1&-1&-1&1\end{pmatrix},
 \qquad H_4^2=4I_4.
\]
Le réseau construit ci-dessous réalise les valeurs positives produites par cette extraction. Sa forme combinatoire est une formalisation ultérieure du registre, avec une origine et une orientation déclarées.

\subsection{Réseau orienté acyclique et déterminants de chemins sans intersection}

Pour un registre positif de longueur \(L\), nous posons
\[
 d_r=K_r/Q>0,\qquad t_0=0,\qquad
 t_j=\sum_{r=1}^jd_r\quad(1\le j\le L).
\]
Ainsi \(t_L=1\). Les \(L\) séparations produisent \(L+1\) points marqués ;
dans le cas dodécaphasique, il y a treize points. La matrice qui les réalise est
\begin{equation}
 C(d)=\begin{pmatrix}1&1&\cdots&1\\t_0&t_1&\cdots&t_L\end{pmatrix}.
 \label{eq:nc-path-matrix}
\end{equation}
La grassmannienne \(\operatorname{Gr}(2,L+1)\) identifie les matrices de rang
deux par changements inversibles des deux lignes. Sa partie positive est formée
des classes dont tous les mineurs ordonnés d’ordre deux
peuvent être strictement positifs.

Nous construisons un réseau orienté de sommets \(a_j,b_j\), pour \(0\le j\le L\), et de puits \(q_j\). Il comporte des arêtes horizontales \(a_{j-1}\to a_j\) et \(b_{j-1}\to b_j\), de poids un ; des arêtes verticales \(a_j\to b_j\), de poids \(d_j\), pour \(j\ge1\) ; et des arêtes \(b_j\to q_j\), de poids un. Les sources ordonnées sont \(b_0,a_0\). Les deux lignes sont dessinées horizontalement, les arêtes verticales descendent et les puits sortent en dessous de la ligne inférieure. Ce dessin est planaire et l'orientation est acyclique. La mesure de frontière \(M_{ij}\) est la somme des produits de poids de tous les chemins allant de la source \(i\) au puits \(q_j\). Le réseau de chemins a deux sources et \(L+1\) puits, soit \(L+3\) connexions extérieures. La matrice qu'il réalise a \(L+1\) colonnes. Le germe plabique introduit ensuite possède précisément \(L+1\) sommets de frontière et constitue une autre présentation de ses coordonnées ; ces deux décomptes de frontière restent explicites.

\begin{theorem}[Mesure de frontière et mineurs du registre]
\label{thm:nc-network}
La mesure de frontière de ce réseau est \(M=C(d)\). Pour
\(0\le i<j\le L\),
\begin{equation}
 \Delta_{ij}(C(d))=t_j-t_i=\sum_{r=i+1}^{j}d_r>0.
 \label{eq:nc-minor-gap}
\end{equation}
Le même nombre est la somme des poids des paires de chemins
à sommets disjoints \(b_0\to q_i\), \(a_0\to q_j\).
\end{theorem}
\begin{proof}
La source inférieure a un chemin unique vers chaque puits, de poids un.
Un chemin de \(a_0\) à \(q_j\) descend exactement une fois, par une
arête d’indice \(1\le r\le j\) ; son poids est \(d_r\). Cela prouve la
formule de la matrice. Le chemin inférieur jusqu’à \(q_i\) occupe tous les
sommets \(b_0,\ldots,b_i\). Par conséquent, un chemin de l’autre source vers
\(q_j\) en est disjoint précisément lorsqu’il descend en
\(i<r\le j\). Sa somme de poids est le membre droit de
\eqref{eq:nc-minor-gap}. Le couple dont les destinations sont échangées
partage toujours un sommet. Le développement du déterminant annule tous les
couples qui se croisent et conserve exactement les couples disjoints. On a ici
vérifié directement, pour le réseau concret, le mécanisme déterminantal
des chemins sans intersection.
\end{proof}

L’identification des mesures de réseaux plans aux cellules de la
grassmannienne non négative fournit leur reconnaissance géométrique ultérieure ;
voir A. Postnikov,
\href{https://arxiv.org/abs/math/0609764}{\emph{Total positivity,
Grassmannians, and networks}}. La preuve précédente donne, en outre, une formule
finie pour chaque mineur sans utiliser de valeur métrologique.

La collection \(d_r>0\), \(\sum_rd_r=1\), a pour dimension \(L-1\). Son application dans la grassmannienne est injective : si les lignes de \(C(d)\) et \(C(d')\) engendrent le même plan, les secondes lignes diffèrent par une transformation affine ; les conditions \(t_0=t'_0=0\) et \(t_L=t'_L=1\) fixent cette transformation. En particulier, pour \(L=12\), on obtient une collection de dimension onze à l'intérieur de la cellule positive de dimension vingt-deux. Le point appartient à la cellule supérieure ; la collection conserve sa dimension propre. Les treize points et les douze séparations ont ainsi des rôles mathématiques distincts.

Autoriser \(d_r\ge0\) conserve la non-négativité. Lorsque \(d_r=0\), les colonnes consécutives \(r-1\) et \(r\) coïncident et forment des éléments parallèles du matroïde. Le rang reste égal à deux car \(t_0=0\) et \(t_L=1\). Les bases sont exactement les paires \(i<j\) qui traversent une séparation positive. Cette dégénérescence produit des strates d'incidence consécutive. Une colonne nulle, qui donnerait une boucle, correspond à une opération distincte de la coïncidence de deux colonnes non nulles.

\subsection{Germe de triangulation et relation d’échange}

Soit \(N=L+1\) et considérons les sommets \(0,\ldots,L\) d’un polygone
convexe dans leur ordre cyclique. À chaque côté ou diagonale \((i,j)\), on assigne
la coordonnée \(A_{ij}=\Delta_{ij}(C(d))\), avec des indices ordonnés.
La triangulation en éventail depuis le sommet zéro a pour diagonales
\((0,j)\), \(2\le j\le L-1\). Ses côtés sont des variables gelées et
ses diagonales des variables mutables. Le carquois s’obtient en orientant
cycliquement les trois côtés de chaque triangle et en annulant les couples de
flèches opposées. C’est un germe de type \(A_{N-3}\) ; pour le registre
dodécaphasique, son type est \(A_{10}\).

Le germe possède une réalisation plabique explicite. On place un sommet
noir à l’intérieur de chaque triangle ; aux sommets du polygone, on place un
sommet blanc s’il est incident à une diagonale, et noir dans le cas contraire.
On relie les centres aux trois sommets de leurs triangles, on retire les
côtés originaux et on contracte les arêtes entre sommets de même couleur.
On ajoute un feuillet de frontière à chaque sommet marqué. Cette construction
reprend l’algorithme de triangulation de Kodama et Williams,
\href{https://people.math.harvard.edu/~williams/papers/KW-ArxivVersion.pdf}
{\emph{KP solitons and total positivity for the Grassmannian}, algorithme 12.1}.
Les étiquettes de coordonnées de ce germe sont évaluées dans les mineurs déjà produits par le réseau à deux lignes. Pour une face mutable, la coordonnée de type \(\mathcal X\) est le quotient des produits des coordonnées \(\mathcal A\) incidentes prescrits par sa colonne du carquois ; dans un quadrilatère, elle sera écrite explicitement comme un birapport. Changer le graphe modifie la coordinatisation de ces mineurs, tout en maintenant leur provenance.

\begin{theorem}[Échange positif de la réalisation du registre]
\label{thm:nc-ptolemy}
Pour quatre sommets \(a<b<c<d\),
\begin{equation}
 A_{ac}A_{bd}=A_{ab}A_{cd}+A_{ad}A_{bc}.
 \label{eq:nc-ptolemy}
\end{equation}
Par conséquent, le changement de diagonale \((a,c)\leftrightarrow(b,d)\)
s’évalue au moyen de
\[
 A_{bd}=\frac{A_{ab}A_{cd}+A_{ad}A_{bc}}{A_{ac}}>0.
\]
Toutes les triangulations fournissent des coordinatisations positives compatibles sur la même réalisation \(C(d)\).
\end{theorem}
\begin{proof}
En substituant \(A_{ij}=t_j-t_i\), l’égalité se réduit à
\[
 (t_c-t_a)(t_d-t_b)
 =(t_b-t_a)(t_d-t_c)+(t_d-t_a)(t_c-t_b).
\]
Le développement des deux membres prouve l’identité. Les différences sont
positives ; chaque flip conserve donc une évaluation positive et le
quotient est défini. La connexion des triangulations par flips
permet de passer de l’une à l’autre au moyen de ces échanges. L’interprétation
en amas de la relation de Plücker est la reconnaissance ultérieure développée
par J. Scott dans
\href{https://arxiv.org/abs/math/0311148}{\emph{Grassmannians and Cluster
Algebras}}. La preuve algébrique de l’évaluation concrète est déjà contenue
dans l’identité montrée.
\end{proof}

\subsection{Raffinement des séparations et naturalité des échanges}

Soit \(B\ge2\). À la profondeur \(n\), chaque séparation \(d_r\) est divisée
en \(B^n\) séparations consécutives égales :
\[
 d^{(n)}_{(r-1)B^n+c}=d_r/B^n,
 \qquad 1\le c\le B^n.
\]
Cela définit \(LB^n+1\) points et une matrice \(C_n\). Si \(\iota_n(j)=Bj\) identifie un point à son descendant de même position, alors
\begin{equation}
 C_{n+1}|_{\iota_n(\{0,\ldots,LB^n\})}=C_n,
 \qquad
 \Delta_{\iota_n(i),\iota_n(j)}(C_{n+1})=\Delta_{ij}(C_n).
 \label{eq:nc-refinement}
\end{equation}
Le réseau raffiné remplace chaque arête verticale de poids \(d\) par une
séquence de \(B\) positions de descente, de poids \(d/B\). La somme
des chemins entre extrémités héritées reste \(d\). La composition des
raffinements conserve cette égalité car les sommes finies se regroupent
de manière associative.

\begin{proposition}[Naturalité sur les extrémités héritées]
\label{prop:nc-inherited}
Tout échange \eqref{eq:nc-ptolemy} entre quatre extrémités héritées commute exactement avec le raffinement. Toute triangulation ancienne s'étend à une triangulation du polygone raffiné en préservant ses triangles intérieurs : les polygones ajoutés entre deux extrémités consécutives sont triangulés à l'extérieur de ces triangles. Les flips entre anciennes diagonales conservent alors leurs relations d'échange.
\end{proposition}
\begin{proof}
La première conclusion résulte du remplacement de chaque mineur par son égalité dans
\eqref{eq:nc-refinement}. Pour la seconde, on conserve les cordes qui
relient les anciennes extrémités consécutives et on triangule chaque poche extérieure
formée par les nouvelles marques. Ces régions ont des intérieurs disjoints.
Les deux triangles incidents à une ancienne diagonale restent identiques,
et donc leur quadrilatère et leur identité de Ptolémée également.
\end{proof}

Cette naturalité spécifie les extrémités sur lesquelles elle agit. Un côté
ancien, lorsqu’il reçoit de nouveaux sommets sur son arc de frontière, devient une
diagonale du polygone plus grand. Une variable anciennement gelée
peut donc être mutable dans le problème étendu. La compatibilité précédente est
une compatibilité exacte du système d’évaluations et des flips
hérités. Une inclusion de motifs de germes ayant le même régime de
coefficients s’obtient en maintenant gelées ces cordes et les nouvelles
diagonales auxiliaires. Cela déclare l’opération qui conserve le germe ;
l’accroissement de profondeur et une mutation demeurent des opérations
distinctes.

\subsection{Action diagonale projective et invariance des birapports}

Si \(\lambda_j>0\), l'action par colonnes \(C\mapsto C\operatorname{diag}(\lambda_0,\ldots,\lambda_L)\) donne
\[
 A_{ij}\longmapsto\lambda_i\lambda_j A_{ij}.
\]
Elle préserve la positivité et toutes les strates définies par l’annulation de
mineurs. Le birapport associé à un quadrilatère est
\begin{equation}
 X_{abcd}=\frac{A_{ab}A_{cd}}{A_{ad}A_{bc}}>0.
 \label{eq:nc-crossratio}
\end{equation}
Les quatre facteurs \(\lambda_a\lambda_b\lambda_c\lambda_d\) se simplifient exactement. Par conséquent, toute action diagonale positive par colonnes laisse \(X_{abcd}\) fixe, même si elle dépend d'un paramètre et a un déterminant égal à un. La modification du poids des colonnes conserve le point projectif de chaque colonne ; leurs birapports expriment cette conservation. Pour que le registre change la forme projective, il doit agir sur les séparations ordonnées.

\subsection{Déformation positive des séparations dirigée par le registre}

Le corpus fixe la représentation des douze positions et son projecteur
orthogonal \(P_3\) sur le secteur tridimensionnel. Avec
\(P_{11}=I_{12}-J_{12}/12\), la direction retrouvée est
\[
 u=P_3P_{11}K=P_3K\in\mathbf1^\perp.
\]
Dans la coordinatisation déclarée, en posant \(r=\sqrt5\), son évaluation exacte est
\begin{align*}
u=(&-495/4-233r/10,\ 403/4+169r/10,\ -403/4-169r/10,\\
   &495/4+233r/10,\ 601/4+1031r/10,\ 203/4+211r/5,\\
   &-203/4-211r/5,\ -601/4-1031r/10,\ 313/4+569r/10,\\
   &162+219r/20,\ -162-219r/20,\ -313/4-569r/10).
\end{align*}
On conserve \(\sum_ru_r=0\) et
\(\|u\|^2=(6638585+2275584\sqrt5)/20>0\).

Pour \(s\in\mathbb R\), nous définissons la déformation des poids par
\begin{equation}
 d_r(s)=\frac{K_re^{su_r}}{\sum_{\ell=1}^{12}K_\ell e^{su_\ell}},
 \qquad t_j(s)=\sum_{r=1}^jd_r(s),\qquad C(s)=C(d(s)).
 \label{eq:nc-gap-flow}
\end{equation}
La formule utilise le registre et sa direction déjà engendrés. C'est une nouvelle réalisation analytique du lecteur ; le paramètre \(s\) mesure sa déformation. Lorsque la sortie interne \(\rho_A=\pi_{\rm HMT}/\varphi_{\rm HMT}^2\) est utilisée, l'écriture \(s=\tau\log\rho_A\) exprime exactement la même collection. Cette reparamétrisation agit après la génération des constantes. L'identification de \(s\) ou \(\tau\) à un temps physique requiert la loi temporelle du lecteur correspondant.

\begin{theorem}[Changement de forme, positivité et compatibilité avec la profondeur]
\label{thm:nc-shape-flow}
La collection \eqref{eq:nc-gap-flow} appartient à \(\operatorname{Gr}_{>0}(2,13)\) pour tout \(s\) réel et conserve \(t_0=0\), \(t_{12}=1\). Pour \(i<j\), soient
\[
 \mu_{ij}(s)=
 \frac{\sum_{r=i+1}^jK_ru_re^{su_r}}{\sum_{r=i+1}^jK_re^{su_r}},
 \qquad
 \mu(s)=\frac{\sum_rK_ru_re^{su_r}}{\sum_rK_re^{su_r}}.
\]
Alors
\begin{equation}
 \frac{d}{ds}\log A_{ij}(s)=\mu_{ij}(s)-\mu(s),
 \quad
 \frac{d}{ds}\log X_{abcd}(s)
 =\mu_{ab}(s)+\mu_{cd}(s)-\mu_{ad}(s)-\mu_{bc}(s).
 \label{eq:nc-crossratio-derivative}
\end{equation}
La déformation commute avec chaque raffinement qui attribue aux successeurs de la séparation \(r\) le poids \(d_r/B\) et la même vitesse \(u_r\).
\end{theorem}
\begin{proof}
Toutes les exponentielles et tous les \(K_r\) sont positifs. La normalisation
fait que la somme des séparations est un ; les mineurs sont des sommes de
séparations positives par \eqref{eq:nc-minor-gap}. Dériver le logarithme de
\[
 A_{ij}(s)=\frac{\sum_{r=i+1}^jK_re^{su_r}}
                      {\sum_rK_re^{su_r}}
\]
donne la première identité. La combinaison de quatre logarithmes dans \eqref{eq:nc-crossratio} annule les quatre moyennes globales avec leurs signes et donne la seconde. Enfin, remplacer \(K_r\) par \(B\) copies de \(K_r/B\) de même \(u_r\) conserve le numérateur et le dénominateur lorsque les successeurs sont regroupés. L'égalité vaut pour tous les paramètres, à toute profondeur et pour toute composition de raffinements.
\end{proof}

Il existe un témoin exact du changement de forme de cette collection. Pour les quatre premières extrémités,
\begin{equation}
 X_{0123}(0)=\frac{1560}{23711},\qquad
 \left.\frac{d}{ds}\log X_{0123}(s)\right|_{s=0}
 =-\frac{309899}{917}-\frac{268596}{4585}\sqrt5<0.
 \label{eq:nc-nontrivial-witness}
\end{equation}
Le calcul utilise uniquement les trois premières séparations \((234,543,140)\) et leurs vitesses déjà récupérées. L'inégalité stricte prouve un mouvement projectif effectif. Les identités de Ptolémée restent exactes au cours de ce mouvement, puisqu'elles s'appliquent aux colonnes \((1,t_j(s))\) pour chaque paramètre. Le registre contrôle ainsi une déformation positive de la forme et, simultanément, un système de coordinatisations demeurant compatible avec elle.

La portée démontrée réunit un réseau planaire de rang deux, sa mesure
explicite, un germe plabique de triangulation, toutes ses mutations par
flips, la naturalité des extrémités héritées et une déformation des
birapports dirigée par \(K\). Chaque opération a un domaine et un
résultat vérifiables. La comparaison de cette réalisation avec la
dynamique temporelle enrichie et avec la branche matérielle utilise leurs
applications de lecture respectives ; aucune identité de coordonnées ne remplace ces applications.
